\chapter{Conclusion}
\label{ch:conclusion}

\section{Synthesis}
The four analytical chapters were designed to triangulate on the same
underlying business question — where does this retailer's profit come
from, and where does it leak away — from three largely independent
angles, and they agree:

\begin{itemize}
  \item \textbf{Descriptively} (\cref{ch:eda}), mean profit and the
        empirical loss rate degrade sharply once line-level discount
        exceeds roughly 20--30\%, with an exact 0\% loss rate at zero
        discount rising to 96.8\% at $\geq$30\% discount.
  \item \textbf{By construction} (\cref{ch:features}), the same threshold
        appears as a clean, monotone table once discount is bucketed:
        every line discounted 50\% or more in this dataset is
        loss-making, without exception.
  \item \textbf{By trained model} (\cref{ch:models}), two independently
        fitted, different-target ensembles (a profit regressor and a
        loss classifier) both rank \texttt{Discount} and its
        bucketed/derived forms as overwhelmingly the most important
        predictor, well ahead of category, region, shipping mode, or
        any customer-history feature.
\end{itemize}
That three different analytical lenses — plain descriptive statistics,
an engineered lookup table, and a fitted non-linear model with formal
held-out evaluation — converge on the same driver is stronger evidence
than any one of them alone, and is the report's central, actionable
finding: \textbf{a discount governance policy (e.g.\ requiring approval
above roughly a 30\% discount) targets the great majority of preventable
loss in this business}, and the loss classifier of \cref{ch:models}
gives a concrete, auditable tool for flagging specific at-risk lines,
catching 93\% of realised 2017 loss dollars on a strict future-year
evaluation.

The secondary line of work, forecasting (\cref{ch:forecasting}),
is a comparatively harder problem on this data — a single 48-month
series is a small sample by time-series standards — but nonetheless
supports a directionally useful conclusion: aggregate sales are on a
clear upward trend (roughly +25\% projected for the next 12 months), with
materially uneven growth by category (Office Supplies accelerating,
Technology flat), information a single store-wide growth assumption
would obscure.

\section{Limitations}
\begin{itemize}
  \item \textbf{Correlation, not causation, on discounting.} The finding
        that deep discounts coincide with losses is robust and
        consistent across every analysis in this report, but none of
        them establishes that \emph{cutting} discounts on the affected
        lines would recover that profit at the same customers' current
        purchase volumes — price elasticity of demand is not modelled,
        and a real governance policy should be piloted (e.g.\ as an
        A/B or staged rollout) rather than deployed store-wide on the
        strength of observational data alone.
  \item \textbf{Small forecasting sample.} $T=48$ monthly observations
        is little more than one seasonal cycle's worth of replication
        (four Decembers), so the 8-fold backtest of \cref{ch:forecasting}
        — while far better than a single split — still rests on
        overlapping, non-independent windows of a short series;
        confidence intervals and model rankings there should be treated
        as indicative rather than tightly precise.
  \item \textbf{No hyperparameter search.} Both the forecasting models
        (SARIMA order, Holt--Winters trend/seasonal type) and the
        supervised models (tree depth, learning rate, number of
        estimators) use fixed, reasonable-default configurations rather
        than a tuned grid/Bayesian search; the reported metrics are
        therefore a floor on achievable performance for each model
        family, not necessarily its best achievable result.
  \item \textbf{Single dataset, single hold-out year.} All supervised
        results are evaluated on one train/test temporal split
        (2014--16 vs.\ 2017); a rolling-origin backtest analogous to
        \cref{ch:forecasting}'s would give a more robust estimate of
        out-of-year generalisation for the supervised models too, at
        the cost of additional retraining.
  \item \textbf{Synthetic/aggregated public dataset.} The Superstore
        dataset is a widely used teaching/demo dataset rather than a
        live production feed; absolute dollar figures and even some
        structural regularities (e.g.\ the very sharp, almost
        deterministic discount-loss threshold) may be cleaner than a
        genuinely messy operational dataset would be, and any
        recommendation here should be re-validated against the target
        business's real data before acting on it.
\end{itemize}

\section{Recommendations}
\begin{enumerate}[label=\arabic*.]
  \item \textbf{Introduce a discount approval gate} around the 30\%
        threshold identified in \cref{tab:discount-bucket,fig:clf-imp},
        prioritised at the sub-category level (Tables, Bookcases first,
        per \cref{tab:subcat-profit}).
  \item \textbf{Operationalise the loss classifier} of \cref{ch:models}
        as a pre-order or end-of-day flag on discounted lines, using the
        catch-rate/false-flag trade-off of \cref{sec:business-framing}
        as the basis for setting the alert threshold to the
        business's actual tolerance for false positives.
  \item \textbf{Plan 2018 category growth unevenly}, not with a single
        store-wide multiplier: increase inventory/staffing allocation
        toward Office Supplies (+32.8\% projected) and hold Technology
        roughly flat (+0.2\% projected), per \cref{tab:cat-forecast}.
  \item \textbf{Pilot before store-wide rollout} of any discount policy
        change, given the correlational nature of the evidence
        (Limitations, above), and re-fit the forecasting and supervised
        models on a rolling basis as new months of data arrive — the
        entire pipeline (\texttt{src/pipeline.py}) is built exactly for
        this kind of periodic re-run.
\end{enumerate}

\section{Future work}
Natural extensions, roughly in order of expected value for effort:
(i) a genuine \emph{uplift}/causal model for discount policy (e.g.\
comparing similar orders that did and did not receive a discount, or a
proper randomised pilot) to move the central finding from correlational
to causal; (ii) hyperparameter tuning (grid or Bayesian search) for both
the forecasting and supervised models, with nested/rolling
cross-validation rather than a single split; (iii) hierarchical
forecasting (reconciling the whole-store and per-category forecasts of
\cref{tab:cat-forecast} to be mutually consistent, e.g.\ via bottom-up or
MinT reconciliation); and (iv) extending the supervised models with
richer product- and customer-level features (e.g.\ text features from
\texttt{Product Name}, or a proper customer lifetime value model) beyond
the RFM segmentation of \cref{ch:features}, which was used here only
descriptively.

\appendix
\chapter{Reproducibility}
\label{app:repro}

\section{Chapter-to-code map}
\begin{table}[H]
  \centering
  \caption{Every figure and number in this report traces to one of these
  notebook/module pairs.}
  \label{tab:repro-map}
  \begin{tabular}{@{}lll@{}}
    \toprule
    Report chapter & Notebook & Core module \\
    \midrule
    \cref{ch:eda}         & \texttt{01\_data\_overview.ipynb}, \texttt{02\_eda.ipynb} & \texttt{src/data\_loader.py} \\
    \cref{ch:features}    & \texttt{03\_feature\_engineering.ipynb} & \texttt{src/features.py} \\
    \cref{ch:forecasting} & \texttt{04\_forecasting.ipynb} & \texttt{src/forecasting.py} \\
    \cref{ch:models}      & \texttt{05\_supervised\_models.ipynb} & \texttt{src/modeling.py} \\
    all (glue)             & \texttt{06\_end\_to\_end.ipynb} & \texttt{src/pipeline.py} \\
    \bottomrule
  \end{tabular}
\end{table}

\section{One-command reproduction}
From \texttt{EY Training/ml\_superstore/}:
\begin{lstlisting}
python -m pip install -r requirements.txt
python -m src.pipeline              # ~15s: writes outputs/REPORT.md
                                     # and outputs/tables/pipeline_*.csv
\end{lstlisting}
Every figure embedded in this report is regenerated by re-executing the
notebooks in numeric order (see \texttt{../README.md}); this document is
then rebuilt with \texttt{pdflatex} against the (unchanged, unless the
underlying data changes) \texttt{outputs/figures/} directory, per
\texttt{report/README.md}.

\section{Notation summary}
\begin{table}[H]
  \centering
  \caption{Symbols used across chapters.}
  \begin{tabular}{@{}ll@{}}
    \toprule
    Symbol & Meaning \\
    \midrule
    $n$, $p$ & sample size, number of features (supervised setting) \\
    $\mathbf x_i$, $y_i$, $\hat y_i$ & feature vector, true target, predicted target for row $i$ \\
    $T$, $y_t$, $\hat y_{t+h}$ & series length, monthly value at $t$, $h$-step forecast (\cref{ch:forecasting}) \\
    $B$, $s$ & backshift operator, seasonal period (12 months) \\
    $\alpha,\beta,\gamma$ & Holt--Winters level/trend/seasonal smoothing parameters \\
    $\lambda$ & Ridge regularisation strength (\cref{eq:ridge}) \\
    $\sigma(\cdot)$ & logistic sigmoid \\
    $\rho$ & Pearson correlation (\cref{ch:eda}) or inter-tree correlation (\cref{ch:models}), disambiguated in context \\
    \bottomrule
  \end{tabular}
\end{table}
